\documentclass[11pt]{article}       % set main text size
\usepackage[letterpaper,            % set paper size to letterpaper. change to a4paper for resumes outside of North America
top=0.5in,                          % specify top page margin
bottom=0.5in,                       % specify bottom page margin
left=0.5in,                         % specify left page margin
right=0.5in]{geometry}              % specify right page margin
                       
\usepackage{XCharter}               % set font
\usepackage[T1]{fontenc}            % output encoding
\usepackage[utf8]{inputenc}         % input encoding
\usepackage{enumitem}               % enable lists for bullet points: itemize and \item
\usepackage[hidelinks]{hyperref}    % format hyperlinks
\usepackage{titlesec}               % enable section title customization
\raggedright                        % disable text justification
\pagestyle{empty}                   % disable page numbering

% ensure PDF output will be all-Unicode and machine-readable
\input{glyphtounicode}
\pdfgentounicode=1

% format section headings: bolding, size, white space above and below
\titleformat{\section}{\bfseries\large}{}{0pt}{}[\vspace{1pt}\titlerule\vspace{-6.5pt}]

% format bullet points: size, white space above and below, white space between bullets
\renewcommand\labelitemi{$\vcenter{\hbox{\small$\bullet$}}$}
\setlist[itemize]{itemsep=-2pt, leftmargin=12pt}

% resume starts here
\begin{document}

% name
\centerline{\Huge Tahar Meijs}

\vspace{5pt}

% contact information
\centerline{\href{mailto:tntmeijs@gmail.com}{tntmeijs@gmail.com} | \href{https://www.linkedin.com/in/tntmeijs/}{linkedin.com/in/tntmeijs} | \href{https://github.com/tntmeijs}{github.com/tntmeijs}}

% experience section
\section*{Experience}
\textbf{\href{https://www.deloitte.com/nl/nl.html}{Deloitte -- Amsterdam, The Netherlands}\hfill September 2020 -- Present} \\
\textit{Specialist lead / manager} \\
\vspace{-9pt}
\begin{itemize}
  \item Built a neobank on AWS with Go, Lambda, and DynamoDB within 12 months that uses a serverless architecture to achieve low average operational costs and support high burst traffic. It also eliminated the need for an expensive Kubernetes cluster with JVM containers.
  \item Implemented a monitoring stack on CloudWatch to track Step Function performance and SQS queue loads.
  \item Designed event-driven serverless architecture using Lambda, EventBridge, and SQS.
  \item Coached and mentored junior engineers, guiding them from being on a performance improvement plan to promotion to the next level within two years.
\end{itemize}

\textit{Senior specialist / senior consultant} \\
\vspace{-9pt}
\begin{itemize}
  \item Led an 8-person engineering team building an event-driven load origination platform, integrating with a dozen external APIs and services, resulting in successful project delivery and 3 team member promotions.
  \item Built a corporate loan origination platform processing 50M+ euros in loans annually, reducing approval time from multiple weeks to days.
  \item Decomposed a distributed monolith application into a dozen domain-specific microservices, improving our team's velocity to deliver new tickets.
\end{itemize}

\textit{Specialist / consultant} \\
\vspace{-9pt}
\begin{itemize}
  \item Built an event-driven retail loan origination platform for a large Dutch bank using Java and AWS, generating 7.5M+ euros in loans each year.
  \item Designed and built a custom debugging tool to resolve bugs in event-driven systems using Kafka, significantly speeding up incident resolution time and improved developer experience.
  \item Built an ETL pipeline processing 100k+ JSON Kafka events per hour, all transformed into SQL table entries, and uploaded to a datalake using FTP for audit purposes.
\end{itemize}

\textit{Analyst} \\
\vspace{-9pt}
\begin{itemize}
  \item Contributed to a digital loan origination solution for a large Dutch bank that reduces the time it takes to apply for a loan from 6 weeks to 15 minutes by automating their offline processes.
\end{itemize}

\textbf{\href{https://www.sumo-digital.com/}{Sumo Digital -- Sheffield, United Kingdom}\hfill September 2019 -- August 2020} \\
\textit{Junior software engineer}\\
\vspace{-9pt}
\begin{itemize}
  \item Built a real-time testing and video streaming dashboard processing over 10.000 test results daily using React, Node.js, and Websockets, reducing the QA feedback cycle from hours to minutes.
  \item Optimised PlayStation 5 resource packing algorithms in C++ by more efficiently allocating data chunks next to each other, reducing the game's overall size by 25\%, allowing more data to be stored on Blu-ray discs.
  \item Developed a cross-platform video recording application for PlayStation 4/5 and PC with minimal overhead by asynchronously capturing the GPU's frame buffer.
\end{itemize}

\textit{Placement Programmer} \\
\vspace{-9pt}
\begin{itemize}
  \item Designed and built a custom, cross-platform, testing framework for PC, PlayStation 4, and PlayStation 5 using C\texttt{++} for PlayStation 5 launch title "\href{https://www.playstation.com/en-us/games/sackboy-a-big-adventure/}{Sackboy: A Big Adventure}".
  \item Built a custom performance profiler for Unreal Engine 4 games with minimal overhead on the game's processes by intelligently batching operations.
\end{itemize}

\vspace{-9pt}

\section*{Notable technical challenges}
\begin{itemize}
  \item Eliminated race conditions affecting roughly 1\% of users in our KYC processes by implementing optimistic concurrency with self-healing, eliminating users stuck in limbo states.
  \item Introduced a custom-built debugger for distributed event-driven systems, improving developer debugging velocity by 50\%. Bugs that would take two hours to resolve could now be resolved within the hour.
  \item Rearchitectured core parts of our application to use transactional writes and centrallised orchestration instead of tightly-coupled peer-to-peer orchestration patterns. This solved nearly all eventual consistency bugs.
  \item Asynchronously captured the GPU's frame buffer with minimal impact on the running processes. The game's logic and rendering took well over 14 milliseconds per frame, leaving less than 3 milliseconds to capture the frame buffer and (asynchronously) start H.264 encoding on NVIDIA's NVENC cores.
\end{itemize}

\vspace{-9pt}

% skills section
\section*{Skills}
\begin{tabular}{@{}l l@{}}
\textbf{Backend technologies} & Go, Java, Spring Boot, Node.js \\
\textbf{Cloud and infrastructure} & AWS (Lambda, Step Functions), Docker, Kubernetes, CDK, Terraform \\
\textbf{Data and messaging} & Kafka, DynamoDB, MongoDB, S3, SQS, EventBridge \\
\textbf{Monitoring and observability} & CloudWatch, DataDog \\
\textbf{Architecture} & Serverless, microservices, event-driven design, domain-driven design \\
\textbf{Compliancy} & PII data masking, GDPR-compliant architecture, audit trails \\
\end{tabular}

\vspace{-9pt}

% education section
\section*{Education}
\textbf{\href{https://www.buas.nl/en/programmes/creative-media-and-game-technologies/programming}{Breda University of Applied Sciences}} -- BSc Creative Media and Game Technologies\hfill 2016 -- 2020

\end{document}
